Seja um grupo \(G\), e denote por \(G(m)\) o subgrupo gerado pelas potências \(m^{\text{th}}\) dos elementos de \(G\). Se \(G(m)\) e \(G(n)\) são comutativos, prove que \(G(\operatorname{gcd}(m,n))\) também é comutativo. (\(\operatorname{gcd}(m,n)\) denota o maior divisor comum de \(m\) e \(n\).)
